% LLNCS LaTeX paper — Springer Proceedings format
% Enhancing Customer Experience and Brand Perception: Hotel Review NLP Pipeline
\documentclass[runningheads]{llncs}

\usepackage[T1]{fontenc}
\usepackage{graphicx}
\usepackage{booktabs}
\usepackage{multirow}
\usepackage{amsmath}
\usepackage{hyperref}
\renewcommand\UrlFont{\color{blue}\rmfamily}
\usepackage{color}

\begin{document}

\title{Enhancing Customer Experience and Brand Perception: Leveraging NLP and Digital Marketing for a Hotel in Porto}

\titlerunning{Hotel Review NLP Pipeline}

\author{
  Andr\'{e} Pires\inst{1}\and
  Manuel Sampaio\inst{1}\and
  Pedro Correia\inst{1}\and
  Pedro Meireles\inst{1}\and
  Mi\l{}osz Koz\l{}uk\inst{1}
}

\authorrunning{A. Pires et al.}

\institute{
  IADE -- Universidade Europeia, Lisboa, Portugal
}

\maketitle

%------------------------------------------------------------------
\begin{abstract}
This paper presents a multi-stage NLP pipeline applied to 999 verified
guest reviews collected from Booking.com for a boutique hotel in Porto,
Portugal. The pipeline encompasses exploratory data analysis, supervised
sentiment classification, unsupervised Latent Dirichlet Allocation (LDA)
topic modelling, K-Means customer segmentation, and rule-based
Aspect-Based Sentiment Analysis (ABSA). The overall average guest score
is 8.24/10 across 38 months of data. The primary operational pain point
identified by ABSA is WiFi \& Check-in, where 39.5\% of aspect mentions
are negative. A binary LinearSVC classifier achieves F1-macro of 0.6471
(5-fold CV: 0.6939) and is proposed for real-time review monitoring.
A fine-tuned XLM-RoBERTa transformer is evaluated as a multilingual
proof-of-concept but underperforms classical baselines at this dataset
size. Seven evidence-backed strategic recommendations with measurable
KPIs are derived, addressing WiFi infrastructure, family guest
experience, negative review response, staff-led marketing, and Iberian
market expansion.

\keywords{Hotel reviews \and Sentiment analysis \and Topic modelling
\and Customer segmentation \and ABSA \and NLP \and Digital marketing}
\end{abstract}

%------------------------------------------------------------------
\section{Introduction}

Online guest reviews represent one of the richest, most candid sources
of intelligence available to hospitality operators. Unlike internal
surveys, they are unprompted, public, and directly influence future
booking decisions. For a hotel experiencing declining satisfaction scores
and stagnating brand perception, systematically mining this corpus offers
both a diagnostic and a strategic compass.

This paper presents the end-to-end findings of a data science pipeline
applied to \textbf{999 verified guest reviews} collected from
Booking.com between March 2023 and April 2026. The pipeline---built
with LangGraph for orchestration and scikit-learn~\cite{scikit-learn}
for NLP modelling---covers exploratory analysis, supervised sentiment
classification, unsupervised topic modelling, K-Means customer
segmentation, and evidence-backed strategic recommendations.

The headline finding is an overall average score of \textbf{8.24/10},
masking significant variation by traveller segment and topic cluster.
The following sections unpack that variation and translate it directly
into marketing and operational actions.

%------------------------------------------------------------------
\section{Data Collection and Preparation}

\subsection{Scraping and Ingestion}

Reviews were collected programmatically from Booking.com using a
custom Selenium-based scraper targeting a single boutique property in
Porto, Portugal. The scraper handles dynamic pagination, modal dialogs,
and deduplication via content-based fingerprinting. Each review record
captures: reviewer name, country, room type, number of nights, date of
stay, traveller type, positive and negative review text, and numerical
score.

\textit{Data ethics note:} Reviewer names and nationalities constitute
personal data under GDPR. Names were collected solely for deduplication
and are not used in any analysis. All data was collected from publicly
accessible pages in compliance with Booking.com's terms of service
\cite{booking2026}.

\subsection{Cleaning and Feature Engineering}

After automated cleaning, \textbf{999 reviews} were retained from an
initial 1,000 (one dropped for empty review text). The pipeline applied:
date parsing to period strings (YYYY-MM); number-of-nights extraction;
concatenation of positive and negative text fields into a unified
\texttt{full\_review\_text} column; language detection with
\texttt{langdetect} and translation of non-English reviews via Google
Translate; and derivation of a three-class sentiment label from the
numerical score (positive: $\geq 8$; neutral: 6--7; negative:
$< 6$).

%------------------------------------------------------------------
\section{Exploratory Data Analysis}

\subsection{Score and Sentiment Distribution}

Scores range from 1.0 to 10.0 on Booking.com's scale, with a mean of
\textbf{8.24} (median 8.0, std 1.68). The distribution is left-skewed:
the majority of reviewers score the hotel 8 or above
(Fig.~\ref{fig:score_dist}). Sentiment labels derived from scores
yield 779 positive (78.0\%), 156 neutral (15.6\%), and 64 negative
(6.4\%) reviews.

\begin{figure}[t]
  \centering
  \includegraphics[width=0.75\textwidth]{fig_score_distribution.png}
  \caption{Score frequency distribution across 999 reviews.}
  \label{fig:score_dist}
\end{figure}

\subsection{Temporal Trends}

Review volume peaks in March 2026 (86 reviews). Stacked-area analysis
of monthly sentiment share (Fig.~\ref{fig:sent_time}) reveals a broadly
stable positive share over 38 months, indicating no systematic
operational deterioration. Short-term spikes in neutral and negative
share correlate with periods of higher review volume, suggesting busy
periods may strain service quality.

\begin{figure}[t]
  \centering
  \includegraphics[width=\textwidth]{fig_sentiment_over_time.png}
  \caption{Monthly sentiment share (stacked area) over 38 months.}
  \label{fig:sent_time}
\end{figure}

\subsection{Word Frequency}

Word cloud analysis (Fig.~\ref{fig:wordclouds}) confirms that positive
reviews centre on \textit{location}, \textit{staff}, \textit{breakfast},
and \textit{friendly}, while negative reviews surface \textit{noise},
\textit{small}, \textit{air conditioning}, and \textit{WiFi} as
recurring pain points.

\begin{figure}[t]
  \centering
  \includegraphics[width=0.48\textwidth]{fig_wordcloud_positive.png}
  \hfill
  \includegraphics[width=0.48\textwidth]{fig_wordcloud_negative.png}
  \caption{Word clouds for positive reviews (left) and negative reviews (right).}
  \label{fig:wordclouds}
\end{figure}

\subsection{Guest Origin and Traveller Mix}

The hotel draws an international audience dominated by the
\textbf{United Kingdom} (238 reviews, 23.8\%), followed by the United
States (145) and Ireland (73). Iberian markets (Spain + Portugal: 89
reviews, 8.9\%) are significantly underrepresented given the hotel's
Porto location---a strategic gap addressed in Section~\ref{sec:recs}.
Traveller type scores are shown in Table~\ref{tab:traveller}.

\begin{table}[t]
\caption{Average Booking.com score by traveller type.}
\label{tab:traveller}
\centering
\begin{tabular}{lrr}
\toprule
Traveller type & Avg score & Reviews \\
\midrule
Couple         & 8.35 & 519 \\
Group          & 8.24 & 130 \\
Solo traveller & 8.13 & 226 \\
Family         & 7.98 & 124 \\
\bottomrule
\end{tabular}
\end{table}

%------------------------------------------------------------------
\section{Customer Segmentation}

K-Means clustering ($k=4$) was applied to 999 reviews using five feature
dimensions: guest score, review word count, number of nights, traveller
type (one-hot encoded), and sentiment (ordinal). Features were
standardised with StandardScaler. The value $k=4$ was selected from an
elbow analysis ($k = 2 \ldots 8$) where inertia reduction flattened
beyond four clusters.

\textit{Limitation:} Clustering was performed at the review level, not
the customer level. A single guest with multiple stays may appear across
different clusters. Future work should aggregate features per unique
reviewer.

The four segments are summarised in Table~\ref{tab:segments}.
\textbf{Family Explorers} carry the highest negative share (9.7\%) and
lowest average score, validating targeted improvements in
Section~\ref{sec:recs}. \textbf{Romantic Couples}---the largest,
highest-scoring segment---represent the strongest base for loyalty
campaigns.

\begin{table}[t]
\caption{K-Means customer segments ($k=4$).}
\label{tab:segments}
\centering
\begin{tabular}{lrrrr}
\toprule
Segment            & Size & Share & Avg score & Neg share \\
\midrule
Romantic Couples   & 518  & 51.9\% & 8.35 & 5.2\% \\
Solo Adventurers   & 226  & 22.6\% & 8.13 & 7.1\% \\
Social Groups      & 131  & 13.1\% & 8.24 & 6.9\% \\
Family Explorers   & 124  & 12.4\% & 7.98 & 9.7\% \\
\bottomrule
\end{tabular}
\end{table}

%------------------------------------------------------------------
\section{Sentiment Analysis}

\subsection{Models and Approach}

Three classifiers were evaluated against a VADER~\cite{vader2014}
lexicon baseline across two task formulations: 3-class (positive /
neutral / negative) and binary (positive vs.\ non-positive). All
supervised models use TF-IDF features (max 8,000 terms, 1--2 n-grams,
sublinear TF scaling) with \texttt{class\_weight=`balanced'} to
compensate for the 78/15.6/6.4\% class imbalance. Evaluation uses
macro-averaged F1. Cross-validation uses 5-fold stratified K-Fold.

Additionally, \texttt{xlm-roberta-base}~\cite{xlmroberta2020} was
fine-tuned for 2 epochs as a multilingual benchmark.

\subsection{Results}

Results are shown in Table~\ref{tab:sentiment} and
Fig.~\ref{fig:model_comparison}. The 3-class macro F1 for LR (0.5293)
reflects a structural \textit{label-signal mismatch}: review text at
the score-7/score-8 boundary is linguistically near-identical, making
the neutral/negative boundary unreliable. Both supervised models
nonetheless outperform the VADER baseline (0.4511).

\begin{table}[t]
\caption{Sentiment model results. CV F1 uses 5-fold stratified K-Fold.}
\label{tab:sentiment}
\centering
\begin{tabular}{llrrr}
\toprule
Model                     & Task   & Accuracy & F1-macro & CV F1-macro \\
\midrule
VADER (baseline)          & 3-class & 0.715 & 0.4511 & --- \\
TF-IDF + LR               & 3-class & 0.740 & 0.5293 & $0.4909 \pm 0.039$ \\
TF-IDF + LinearSVC        & 3-class & 0.800 & 0.5344 & $0.4436 \pm 0.043$ \\
TF-IDF + LinearSVC        & Binary  & 0.750 & \textbf{0.6471} & $0.6939 \pm 0.011$ \\
XLM-RoBERTa (fine-tuned)  & 3-class & 0.780 & 0.2921 & --- \\
\bottomrule
\end{tabular}
\end{table}

\begin{figure}[t]
  \centering
  \includegraphics[width=0.80\textwidth]{fig_sentiment_model_comparison.png}
  \caption{F1-macro comparison across all sentiment models.}
  \label{fig:model_comparison}
\end{figure}

\subsection{XLM-RoBERTa Discussion}

XLM-RoBERTa achieved F1-macro 0.2921---\textbf{below the VADER
baseline}. This is expected: transformer models require substantially
more labelled data than the ${\sim}800$ training samples available here.
With only 2 fine-tuning epochs, the model has insufficient exposure to
hotel-review vocabulary. This experiment should be interpreted as a
\textit{proof-of-concept for the multilingual upgrade path}, not a
performance improvement. At the current dataset size, TF-IDF +
LinearSVC remains the recommended production approach.

\subsection{Operational Application}

The binary classifier (F1-macro 0.6471, CV 0.6939) is the most
deployable model. Recommended use: \textbf{auto-flag incoming
non-positive reviews (score $< 8$) within minutes of publication},
triggering a 24-hour management response SLA.

%------------------------------------------------------------------
\section{Topic Modelling}

Latent Dirichlet Allocation (LDA)~\cite{blei2003} was applied to all
999 review texts. The vocabulary was built with CountVectorizer (max
3,000 features, 1--2 n-grams, min\_df=3, max\_df=90\%), yielding 1,289
unique terms after hotel-domain and multilingual stop-word filtering.
Six topics were selected after comparing perplexity scores across
$k \in \{4, 6, 8\}$; $k=6$ produced the most interpretable separation
while avoiding over-fragmentation at $k=8$. Perplexity: \textbf{844.65}.

Discovered topics are shown in Table~\ref{tab:topics} and
Fig.~\ref{fig:topic_keywords}. \textbf{Location \& Accessibility} is
the dominant theme by volume (362 reviews). \textbf{Staff \& Service}
records the highest average score (8.64/10) and is the hotel's primary
reputation differentiator. \textbf{Room Comfort \& Cleanliness} records
the lowest score (7.36/10). The \textit{Check-in \& WiFi} topic (7
reviews) is statistically thin and should be interpreted with caution;
the ABSA findings in Section~\ref{sec:absa} provide more robust
evidence for this pain point.

\begin{table}[t]
\caption{LDA topics discovered ($k=6$). Topics sorted by review count.}
\label{tab:topics}
\centering
\begin{tabular}{lrrp{4.8cm}}
\toprule
Topic label                  & Reviews & Avg score & Top keywords \\
\midrule
Location \& Accessibility    & 362 & 8.16 & location, breakfast, staff, close, bathroom \\
Staff \& Service             & 280 & 8.64 & staff, location, friendly, helpful, clean \\
Food \& Breakfast            & 219 & 8.32 & location, breakfast, staff, excellent, perfect \\
Room Comfort \& Cleanliness  & 112 & 7.36 & parking, location, check, old, bed \\
Rooftop Bar \& Ambiance      & 19  & 8.47 & pillows, design, comfort, doors, extremely \\
Check-in \& WiFi             & 7   & 7.43 & hear, people, dirty, charged, buy \\
\bottomrule
\end{tabular}
\end{table}

\begin{figure}[t]
  \centering
  \includegraphics[width=\textwidth]{fig_topic_keywords.png}
  \caption{Top keywords per LDA topic (relative weight).}
  \label{fig:topic_keywords}
\end{figure}

%------------------------------------------------------------------
\section{Aspect-Based Sentiment Analysis}
\label{sec:absa}

ABSA was applied using a rule-based approach: for each of eight
hotel-domain aspects, keyword occurrences were located in review text
and a $\pm 12$-word context window was scored with VADER. A total of
\textbf{2,156 deduplicated aspect mentions} were extracted across all
999 reviews.

Results are shown in Fig.~\ref{fig:absa_neg} and
Table~\ref{tab:absa}. \textbf{WiFi \& Check-in} is the highest pain
point with \textbf{39.5\% negative mentions} across 114 reviews.
\textbf{Noise} follows at 32.3\%. \textbf{Location} generates the most
discussion (511 mentions) at only 7.2\% negative---the hotel's clearest
marketing asset. The traveller $\times$ aspect heatmap
(Fig.~\ref{fig:absa_heat}) reveals that Family guests show elevated
negative rates for Room and Noise aspects relative to Couples and Solo
travellers.

\begin{table}[t]
\caption{ABSA results sorted by negative mention share.}
\label{tab:absa}
\centering
\begin{tabular}{lrrr}
\toprule
Aspect           & Mentions & Negative & Neg \% \\
\midrule
WiFi \& Check-in & 114 & 45 & 39.5\% \\
Noise            & 99  & 32 & 32.3\% \\
Room             & 444 & 96 & 21.6\% \\
Breakfast        & 314 & 48 & 15.3\% \\
Value            & 74  & 10 & 13.5\% \\
Cleanliness      & 176 & 19 & 10.8\% \\
Staff \& Service & 424 & 37 & 8.7\%  \\
Location         & 511 & 37 & 7.2\%  \\
\bottomrule
\end{tabular}
\end{table}

\begin{figure}[t]
  \centering
  \includegraphics[width=0.80\textwidth]{fig_absa_neg_by_aspect.png}
  \caption{Negative mention share (\%) by aspect. Dashed line shows mean.}
  \label{fig:absa_neg}
\end{figure}

\begin{figure}[t]
  \centering
  \includegraphics[width=\textwidth]{fig_absa_heatmap.png}
  \caption{Negative mention \% by traveller type $\times$ aspect (heatmap).}
  \label{fig:absa_heat}
\end{figure}

%------------------------------------------------------------------
\section{Strategic Recommendations}
\label{sec:recs}

Seven evidence-backed recommendations are derived directly from the
pipeline artifacts. Table~\ref{tab:recs} provides a decision-table
summary with owners and timelines.

\begin{table}[t]
\caption{Strategic decision table. All KPIs measured within stated timeline.}
\label{tab:recs}
\centering
\small
\begin{tabular}{lp{3.5cm}p{3.0cm}l}
\toprule
ID & Finding & KPI & Timeline \\
\midrule
R1 & WiFi \& Check-in: 39.5\% neg ABSA & Topic avg $\geq 8.50$ & 0--6 mo \\
R2 & Family score 7.98 (lowest) & Family avg $\geq 8.20$ & 0--9 mo \\
R4 & 6.4\% negative, unresponded & Neg share $\leq 4\%$ & 0--3 mo \\
R3 & Staff \& Service avg 8.64 & $+10\%$ direct bookings & 3--12 mo \\
R5 & Iberian share 8.9\% in Porto & Iberian share $\geq 15\%$ & 6--18 mo \\
R6 & Room Comfort avg 7.36 & Topic avg $\geq 8.40$ & 3--9 mo \\
R7 & Value aspect 13.5\% neg & Value avg $\geq 8.60$ & 3--6 mo \\
\bottomrule
\end{tabular}
\end{table}

Key high-priority recommendations:

\textbf{R1 --- WiFi \& Check-in:} Audit signal strength across all
floors; install repeaters in weak zones; introduce digital key/PIN
access to eliminate late-night lockout issues.

\textbf{R2 --- Family Experience:} Offer interconnecting rooms and cot
availability; add child-oriented breakfast items; distribute a printed
city guide for families.

\textbf{R4 --- Negative Review Response:} Deploy the binary sentiment
classifier to auto-flag non-positive reviews within minutes of
publication, triggering a 24-hour management response SLA.

\textbf{R3 --- Staff Marketing:} Feature authentic guest quotes about
staff in OTA listings; launch a `Staff Spotlight' social media series;
nominate consistently praised staff for hospitality awards.

%------------------------------------------------------------------
\section{Conclusions}

This analysis of 999 guest reviews delivers a clear, data-driven picture
of hotel performance. An average score of 8.24/10 and 78.0\% positive
sentiment confirm that the hotel's core proposition---friendly staff,
unique vintage character, and central Porto location---resonates strongly
with guests.

Three findings demand immediate operational attention: (1) WiFi and
check-in friction anchors the lowest-scoring ABSA aspect (39.5\%
negative), fixable with targeted infrastructure investment; (2) Family
Explorers are the lowest-scoring segment (7.98/10), representing an
underserved market; and (3) 6.4\% of reviews are negative---a
manageable number that the binary classifier (F1-macro 0.6471) can
surface in real time.

The hotel's greatest strength---staff quality---should be leveraged
aggressively in OTA copy, social media, and PR to reinforce brand
differentiation and justify premium pricing. Medium-term growth lies in
Iberian market expansion (currently 8.9\% of reviews), converting
Romantic Couples into repeat bookers via a loyalty programme, and
standardising room comfort across all room types.

\subsubsection*{Limitations}
\begin{itemize}
  \item \textit{Label-signal mismatch:} Score-derived sentiment labels
    conflate linguistically similar positive/neutral reviews, capping
    3-class F1-macro at ${\sim}0.53$.
  \item \textit{Multilingual corpus:} Non-English reviews introduce
    noise into LDA keyword lists despite stop-word filtering.
  \item \textit{Single property:} Findings are specific to one hotel;
    cross-property benchmarking would strengthen recommendation
    confidence.
  \item \textit{Review-level clustering:} K-Means segments reviews, not
    customers; future work should aggregate per unique reviewer.
\end{itemize}

%------------------------------------------------------------------
\begin{thebibliography}{8}
\sloppy

\bibitem{scikit-learn}
Pedregosa, F., et al.: Scikit-learn: Machine learning in Python.
\textit{Journal of Machine Learning Research} \textbf{12}, 2825--2830
(2011)

\bibitem{vader2014}
Hutto, C.J., Gilbert, E.: VADER: A parsimonious rule-based model for
sentiment analysis of social media text. In: \textit{Proceedings of
ICWSM 2014} (2014)

\bibitem{blei2003}
Blei, D.M., Ng, A.Y., Jordan, M.I.: Latent Dirichlet Allocation.
\textit{Journal of Machine Learning Research} \textbf{3}, 993--1022
(2003)

\bibitem{xlmroberta2020}
Conneau, A., et al.: Unsupervised cross-lingual representation learning
at scale. In: \textit{Proceedings of ACL 2020}, pp. 8440--8451 (2020)

\bibitem{booking2026}
Booking.com: Guest review data collected from publicly accessible review
pages for a boutique hotel in Porto, Portugal. Retrieved 2026.
\url{https://www.booking.com} (accessed 2026)

\end{thebibliography}

\end{document}
